\chapter{The Extension Problem}

We would like to understand
how a symmetric graph $(G, \Gamma)$ can be ``unfolded''
into a larger imprimitive symmetric graph $(\widetilde{G}, \widetilde{\Gamma})$
admitting the original graph as a quotient.
In particular we'd like to understand how pairs of graphs
$(\Gamma, \widetilde{\Gamma})$ where $\widetilde{\Gamma}$
is an ``extension'' of $\Gamma$ are related combinatorially.
That is, we'd like to be able to describe the the structure
of the graph $\widetilde{\Gamma}$
in terms of the structure of the graph $\Gamma$.

The quotient of a symmetric graph contains considerably less information
than the original graph.
Gardiner and Praeger \cite{GPL1993} observed that
some of the information that is lost
may be recovered from the induced bipartite graph 
between adjacent blocks of the partition,
and from a combinatorial design induced on the blocks themselves.

In this chapter we describe Gardiner and Praeger's observations
as well as some of the questions is raises.

\section{Induced Bipartite graph}

% Introductory remarks

Let $\Gamma$ be an imprimitive $G$-symmetric graph,
with $\cB$ a non-trivial $G$-invariant partition of the vertices.
As always, we assume that the quotient $\Gamma_\cB$ has valency at least one,
so the blocks of $\cB$ are independent sets.
For any arc $(B,C)$ of the quotient,
the subgraph of $\Gamma$ induced by $B \cup C$ must be bipartite --
possibly containing some isolated vertices.
If we restrict ourselves to the subgraph induced by
$(\Gamma(C) \cap B) \cup (\Gamma(B) \cap C)$
then we obtain a bipartite graph with no isolated vertices.
We call this graph the {\it induced bipartite graph} of $(B,C)$,
and denote it by $\Gamma[B,C]$.

\begin{proposition}
The induced bipartite graph $\Gamma[B,C]$ is $G_{B \cup C}$ symmetric.
\begin{proof}
Let $(\alpha,\beta)$ and $(\gamma,\delta)$ be any two arcs of $\Gamma[B,C]$.
Without loss of generality we may assume that $\alpha \in B$ and $\beta \in C$.
By the $G$-symmetry of $\Gamma$ we can find $g \in G$ such that
$(\alpha,\beta)^g = (\gamma,\delta)$.
Since $\Gamma[B,C]$ is bipartite, either $\gamma \in B$ and $\delta \in C$ 
in which case $B^g = B$ and $C^g = C$,
or $\gamma \in C$, $\delta \in B$, in which case $B^g = C$ and $C^g = B$.
Either way, $g \in G_{B \cup C}$ and the result follows.
\end{proof}
\end{proposition}


\begin{proposition}
For any two arcs in the quotient, the induced bipartite graphs are isomorphic.
\begin{proof}
Let $(B,C)$ and $(D,E)$ be any two arcs of $\Gamma_\cB$.
Since $\Gamma_\cB$ is $G$-symmetric we can find some $g \in G$
such that $(B,C)^g = (D,E)$. Clearly
$$(\Gamma(C) \cap B)^g = (\Gamma(C^g) \cap B^g) = (\Gamma(E) \cap D),$$
and
$$(\Gamma(B) \cap C)^g = (\Gamma(B^g) \cap C^g) = (\Gamma(D) \cap E).$$
So $g$ induces a bijection
$$(\Gamma(C) \cap B) \cup (\Gamma(B) \cap C)
\to (\Gamma(E) \cap D) \cup (\Gamma(D) \cap E).$$
Since $g$ is an automorphism, adjacency is preserved and we have an isomorphism
from $\Gamma[B,C]$ to $\Gamma[D,E]$.
\end{proof}
\end{proposition}

%Note that if $(B,C)$ and $(D,E)$ are any two arcs of $\Gamma_\cB$
%then $G_{B \cup C}$ is isomorphic to $G_{D \cup E}$
%though they are not actually {\it the same} group.
%They are different isomorphic subgroups of $G$.
%Let $\cA$ denote the isomorphism class of groups
%to which $G_{B \cup C}$ and $G_{D \cup E}$ both belong.
%The two grphs $\Gamma[B,C]$ and $\Gamma[D,E]$
%are $\cA$-equivalent, but not $\cA$-isomorphic.


When passing from a symmetric graph to its quotient
we preserve the adjacency structure of the blocks,
and discard the exact details of which vertices are connected to which others.
Some of the information that has been lost is recoverable from the induced bipartite graph.
But the induced bipartite graph reveals only the {\it local} connectivity.
To reconstruct the original graph from its quotient we need to know
how these induced bipartite graphs are ``glued together''.
Gardiner and Praeger observed that some of this global information about
the way the induced bipartite graphs fit together is captured in a combinatorial design
induced on each of the blocks.
Before we can describe this design we need a little background.



\section{Combinatorial Designs}

% Need to talk about connection between design theory and geometry.

%Connection with group theory:
%Many interesting groups, including a number of
%the sporadic simple groups arise as automorphism
%groups of $t$-designs.

\begin{definition}
An {\it incidence structure} is a triple $(P,B,I)$
where $P$, and $B$ are sets, usually refered to as
{\it points} and {\it blocks} respectively,
and $I$ is an {\it incidence relation} $I \subseteq P \times B$.
\end{definition}

It is often convenient to visualize an incidence structure
as a bipartite graph. 
Take the points as vertices of one side of
the bipartition and the blocks as vertices of the other.
Draw an edge between the vertex corresponding to a point $p$
and the vertex corresponding to a block $b$ if and only if
$(p,b) \in I$.

We will be interested in incidence structures
satisfying strong regularity and symmetry conditions.

\begin{definition}
A $(v,k,\lambda)$-design is an incidence structure such that:
\begin{enumerate}
\item There are $v$ points in total
\item Each block is incident with exactly $k$ points
\item Each point is incident with exactly $\lambda$ blocks
\end{enumerate}
\end{definition}

When visualizing an incidence structure as a bipartite graph,
the extra regularity conditions of a design correspond 
to the condition that 
any two vertices in the same bipartite half of the graph
must have the same valency.

Note that our definition here of a ``design'' 
corresponds to what is usually refered to as a $1$-design 
or ``tactical configuration'' in the literature.
There is a more general definition of a $t$-design
of which the definition of a $1$-design is a special case.
We will not be needing this more general definition.

% Connection with geometry?

%As an example, the points and lines of the projective plane 
%over a field of $q$ elements, for $q$ a prime power, 
%form a $(q^2 + q + 1,q + 1,q + 1)$-design. 
%The points and lines of the affine plane over a field with $q$ elements 
%form a $(q^2,q,q)$ design. 
%In a very loose sense designs, particularly $2$-designs,
%may be viewed as generalized finite geometries.


The incident point-block pairs of a design are usually refered to as {\it flags}.
For a design $\cD$ we will use the notation
$P_\cD$, $B_\cD$ and $F_\cD$ to denote the points,
blocks and flags of $\cD$ respectively.

Define the {\it trace} of a block $b$ to be the set
$T(b) = \{ p \in P : (p,b) \in I \}$.
Simmilarly define the {\it trace} of a point $p$
to be the set $T(p) = \{b \in B : (p,b) \in I \}$.

\begin{proposition}
In a design, the number of blocks with the same trace is a constant 
that is independent of the choice of block
\end{proposition}

We denote this constant by $m$ and call it the 
{\it multiplicity} of the design.
If distinct blocks have distinct traces,
then we may identify the blocks with their traces,
and take $B$ to be a subset of the power set of $P$.
Otherwise we say that the design contains {\it repeated blocks}.

\begin{proposition}
If a design contains repeated blocks,
then the incidence structure obtained by identifying
blocks with the same trace is again a design
\end{proposition}

\begin{definition}
if $\cD = (P_\cD, B_\cD, I)$ is a $1$-design then the {\it dual} 
of $\cD$ is the design $ \cD^{*} = (B_\cD, P_\cD, I^{*})$
where $(b,p) \in I^{*}$ if and only if $(p,b) \in I$.
\end{definition}

\begin{proposition}
The dual of a $1$-design is a $1$-design.
\end{proposition}



\section{Flag Transitive Designs}

An automorphism of a design $\cD$ is a pair $(\eta, \mu)$ of bijections
$\eta : P_\cD \to P_\cD$ and $\mu : B_\cD \to B_\cD$
with the property that $(p,b)$ is a flag of $\cD$
if and only if $(\eta(p), \mu(b))$ is a flag.
The automorphisms of a design form a group.

As with sets and graphs, we say that a group $G$ acts on a design $\cD$
as a group of automorphisms if there is some homomorphism $\rho : G \to \Aut(\cD)$.
We don't require this homomorphism to be injective, that is we don't
require that $G$ acts {\it faithfully} on $\cD$.

If the group $G$ acts on the design $\cD$ as a group of automorphisms
then we get three induced permutation groups, namely:

\begin{enumerate}
\item $(G, P_\cD)$ -- The induced permutation group on the points.
\item $(G, B_\cD)$ -- The induced permutation group on the blocks.
\item $(G, F_\cD)$ -- The induced permutation group on the flags.
\end{enumerate}


We will be interested in highly symmetric designs.
In particular, we will be interested in pairs $(G, \cD)$
where $G$ acts on $\cD$ in such a way that the induced
permutation group on the flags $(G, F_\cD)$ is transitive.
We will either call the pair $(G, \cD)$ a flag-transitive design,
or we call the design $\cD$ a $G$-flag transitive design.
Many interesting groups arise naturally as flag-transitive
automorphism groups of designs, including a number
of the sporadic simple groups.

\begin{proposition}
If $(G,\cD)$ is a flag-transitive design
then the induced permutation groups $(G,P_\cD)$ and $(G,B_\cD)$
are both transitive
\end{proposition}

%Suppose that $\cD$ is a $G$-flag transitive design, and $(p,b)$ is some flag.
%If $H$ is the stabilizer of $p$, then $(G,P_\cD)$ is a transitive permutation group
%which is permutationally isomorphic to $(G,\cos(H))$.
%If $S$ is the stabilizer of $b$, then $(G,B_\cD)$ is a transitive permutation group
%which is permutationally isomorphic to $(G,\cos(S))$.

\begin{proposition}
If $(G,\cD)$ is a flag-transitive design and $H_1 < G$
is the stabilizer of some point $p$ then the
induced permutation group $(H_1, T(p))$ is transitive.
Simmilarly if $H_2 < G$ is the stabilizer of some block $b$
then the induced permutation group $(H_2, T(b))$ is transitive.
\end{proposition}

%Let $R = S \cap H$, then $R$ stabilizes the flag $(p,b)$.
%Let $B(p)$ denote the blocks of $\cD$ incident with $p$
%and let $P(b)$ denote the points of $\cD$ incident with $b$.
%The action of $H$ on $B(p)$ is permutationally isomorphic to $(H,\cos(R))$
%and the action of $S$ on $P(b)$ is permutationally isomorphic to $(S, \cos(R))$.

If $\cD_1$ and $\cD_2$ are two $G$-flag transitive designs, 
then a {\it $G$-design homomorphism} from $\cD_1$ to $\cD_2$
is a pair of maps $\rho = (\rho_P, \rho_B)$
such that:

\begin {enumerate}
\item $\rho_P : P_{\cD_1} \to P_{\cD_2}$ induces a permutation homomorphism between 
$(G,P_{\cD_1})$ and $(G,P_{\cD_2})$.
\item $\rho_B : B_{\cD_1} \to B_{\cD_2}$ induces a permutation homomorphism between 
$(G,B_{\cD_1})$ and $(G,B_{\cD_2})$.
\item $(p,b) \in I_1$ if and only if $(\rho_P(p), \rho_B(b)) \in I_2$.
\end{enumerate}

When both $\rho_P$ and $\rho_B$ are bijections we have a {\it $G$-design isomorphisms}.
When confusion is unlikely to result, 
we will sometimes write $\rho(p)$ instead of $\rho_P(p)$ when $p$ is a point
and $\rho(b)$ instead of $\rho_B(b)$ when $b$ is a block.

\begin{proposition}
If $\rho$ is a $G$-design homomorphism from $\cD_1$ to $\cD_2$
then $\rho$ induces a $G$-permutation homomorphism from $F_{\cD_1}$ to $F_{\cD_2}$.
\end{proposition}

%For a fixed $G$ the set of all $G$-flag transitive designs
%with $G$-design homomorphisms between them form a category which we shall denote by
%$G$-Flag.
\section{$G$-symmetric designs}

In this section we show how $G$-symmetric designs
may be viewed as a special kind of flag-transitive design.
I'm calling them symmetric designs, but I think
they should probably be called square designs.

\begin{definition}
A $G$-flag transitive design $\cD$ is {\it self-dual}
if there exists a $G$-isomorphism $\rho = (\rho_P, \rho_B)$
between $\cD$ and $\cD^{*}$.
The $G$-isomorphism $\rho$ is called a {\it duality} of $\cD$.
\end{definition}

\begin{definition}
A {\it $G$-symmetric design} is self-dual $G$-flag transitive design $\cD$
which admits a dulaity $\rho = (\rho_P, \rho_B)$ with the property that
$\rho_B \circ \rho_P = id_P$ and $\rho_P \circ \rho_B = id_B$.
In this case the $G$-isomorphism $\rho$ is called a {\it polarity}
of $\cD$.
\end{definition}

Observe that if $\cD$ is a $G$-symmetric design with the property
that $(p, \rho(p))$ is a a flag for some $p \in P_\cD$,
then since $\rho$ is a $G$-isomorphism it follows
that the point-stabilizer of $\cD$ is isomorphic to the flag-stabilizer of $\cD$.
Thus each point is incident with exactly one block
and vice versa.
We consider such $G$-symmetric designs to be ``degenerate'',
and unless an explicit statement to the contrary is given shall
take the expression ``$G$-symmetric design''
to mean ``non-degenerate $G$-symmetric design''.

In a sense that will become clear a little later,
these ``degenerate'' $G$-symmetric
designs correspond to the ``degenerate'' orbital graph $\orb_\Delta(G,\Omega)$
which could be formed from the permutation group $(G, \Omega)$
by taking $\Delta$ to be the diagonal orbit,
and also to the ``degenerate'' Sabidussi graph $\sab(G,H,a)$
which could be formed by taking $a \in H$.
\hfill \break

We shall see that each $G$-symmetric graph
gives rise in a natural way to a $G$-symmetric design
and conversely a $G$-symmetric design together
with a ``marked'' polarity give rise to a $G$-symmetric graph.
We shall also see that in some cases, by choosing a different
polarity it is possible to construct two non-isomorphic $G$-symmetric
graphs from the same $G$-symmetric design.
For any graph $\Gamma$, let $N \Gamma = \{ \Gamma(v) : v \in V \Gamma \}$
denote the set of {\it neighbourhoods} of $\Gamma$.

\begin{proposition}
If \/ $\Gamma$ is a $G$-symmetric graph,
then the incidence structure $\cD(\Gamma) = (V \Gamma, N \Gamma,I)$
where $(v,n) \in I$ if and only if $v \in n$ is a $G$-symmetric design
\begin{proof}
The $G$-arc transitivity of $\Gamma$ is sufficient to ensure
that $\cD(\Gamma)$ is a $G$-flag transitive design.
To see that it is in fact
a $G$-symmetric design we must exhibit a polarity.

Let $\rho_P : P \to B$ be given by $v \mapsto \Gamma(v)$.
The map $\rho_P$ is clearly bijective, and since $G$ acts on $\Gamma$
as a group of automorphisms we have:
\begin{eqnarray*}
\rho_P(v^g) & = & \Gamma(v^g) \\
            & = & \Gamma(v)^g \\
	                & = & \rho_P(v)^g.
			\end{eqnarray*}
			So $\rho_P$ induces a permutation isomorphism.
			Now, take $\rho_B = \rho_P^{-1}$.
			We must check that the pair $\rho = (\rho_P, \rho_B)$ preserve
			the incidence structure of the design.

Suppose that $(v,\Gamma(w))$ is a flag of $\cD$, so $v \in \Gamma(w)$.
Since $\Gamma$ is a simple graph it follows immediately that $w \in \Gamma(v)$
and so $(w, \Gamma(v))$ is also a flag of $\cD$.
That is $(\Gamma(v), w)$ is a flag of $\cD^{*}$.
But
\begin{eqnarray*}
\rho((v, \Gamma(w))) & = & (\rho_P(v), \rho_B(\Gamma(w))) \\
                     & = & (\Gamma(v), w).
		     \end{eqnarray*}
		     So we are done.

		     \end{proof}
		     \end{proposition}

\begin{proposition}
If $\cD$ is a (non-degenerate) $G$-symmetric design with ``marked'' polarity $\rho$
then the graph $\Gamma(\cD, \rho)$ with vertex set $P_\cD$
and arc set $\{ (p,q) : (q,\rho_P(p)) \in I_\cD \}$ is $G$-symmetric.
\begin{proof}
We must first check that $\Gamma(\cD, \rho)$ is well-defined.
The non-degenerateness of $\cD$ ensures that there are no loops.
If $(p,q)$ is an arc of $\Gamma(\cD, \rho)$
then $(q, \rho_P(p))$ is a flag of $\cD$.
Since $\rho$ is an isomorphism,
if $(q, \rho_P(p))$ is a flag of $\cD$
then $(\rho_P(q), \rho_B \circ \rho_P(p)) = (\rho(q), p)$
is a flag of $\cD^{*}$.
It follows that $(p,\rho(q))$ is a flag of $\cD$,
and thus $(q,p)$ is an arc of $\Gamma(\cD, \rho)$.
So $\Gamma(\cD, \rho)$ is simple.


By proposition [?] $G$ acts transitively on the points $\cD$
so $\Gamma(\cD, \rho)$ is $G$-vertex transitive.
Suppose that $(p,q_1)$ and $(p,q_2)$ are two distinct arcs of
$\Gamma(\cD, \rho)$. Then we have $q_1, q_2 \in T(\rho(p))$.
By proposition [?]
the stabilizer of $\rho(p)$ acts transitively on $T(\rho(p))$.
Thus we can find some $g \in G$ which fixes $\rho(p)$ and
carries $q_1$ to $q_2$.
Since $\rho$ is a $G$-isomorphism, if $g$ fixes $\rho(p)$
it also fixes $p$, thus $g$ carries the arc $(p,q_1)$
to the arc $(p,q_2)$. So $\Gamma(\cD, \rho)$ is $G$-locally
transitive. The result follows.

\end{proof}
\end{proposition}

\begin{proposition}
For any design $\cD$ and any polarity $\rho$ the design $\cD(\Gamma(\cD,\rho))$
is isomorphic to $\cD$
\begin{proof}
later
\end{proof}
\end{proposition}

%\begin{definition}
%If $\cD_1$ and $\cD_2$ are two $G$-symmetric designs
%then a $G$-design homomorphism from $\cD_1$ to $\cD_2$
%is a pair of permutation homomorphisms $\eta_P : P_1 \to P_2$ and $\eta_B : B_1 \to B_2$
%satisfying $(p,b) \in I_1$ implies $(\eta_P(p),\eta_B(b)) \in I_2$.
%\end{definition}

%\begin{proposition}
%If $\cD_1$ and $\cD_2$ are two $G$-symmetric designs
%then a $G$-design homomorphism from $\cD_1$ to $\cD_2$
%induces a a $G$-permutation homomorphism from the flags of $\cD_1$
%to the flags of $\cD_2$
%\end{proposition}

%What I want to say is that $G$-symmetric designs
%together with $G$-design homomorphisms form a category
%(well, kind of more of a poset really).
%I also want to say that the category $G$-Graph
%``projects down'' onto the category $G$-Design.
%$F : G$-Graph $\to G$-Design.
%If there's a $G$-graph homomorphism from $x$ to $y$
%then there's a $G$-design homomorphism from $F(x)$ to $F(y)$



\section{``cross-section'' of a graph homomorphism}

Suppose that $\Gamma$ is an imprimitive $G$-symmetric graph
with $\cB$ a non-trivial $G$-invariant partition of the vertices.
Assume further that the quotient $\Gamma_\cB$ has valency
at least one so that the blocks of $\cB$ are independent sets.

For any vertex $\alpha$ of $\Gamma$, let $B(\alpha)$ denote
the block of $\cB$ containing $\alpha$.
Let $\Gamma(\alpha) = \{ \beta \in V \Gamma : (\alpha, \beta) \in A \Gamma \}$
denote the neighbours of $\alpha$ in $\Gamma$.
Let $\Gamma_\cB(B) = \{ C \in \cB : (B,C) \in A \Gamma_\cB \}$
denote the neighbours of $B$ in the quotient.
For any $\alpha \in B$, let
$\Gamma_\cB(\alpha) = \{ B(\beta) : \beta \in \Gamma(\alpha) \}$
denote the blocks of $\cB$ containing the neighbours of $\alpha$ in $\Gamma$.

Construct a design with point set $B$, block set $\Gamma_\cB(B)$
and an incidence relation $ I \subseteq B \times \Gamma_\cB(B)$
defined by $(\alpha, C) \in I$ if and only if $C \in \Gamma_\cB(\alpha)$.

\begin{proposition}
The incidence structure $\cD(B) = (B, \Gamma_\cB(B), I)$ is a $H$-flag transitive
$(v,k,\lambda)$ design, where:
\begin{eqnarray*}
v & := &  |B| \\
k & := &  |\Gamma(C) \cap B|, \\
\lambda  & := &  |\Gamma_\cB(\alpha)| \\
\end{eqnarray*}
and $H$ is the stabilizer of the block $B \in \cB$.
\end{proposition}
	      
The design $\cD(B)$ gives, in a sense, a ``cross-section''
of the homomorphism $\Gamma \mapsto \widetilde{\Gamma}$.

\section{Reconstruction problem}

Gardiner and Praeger observed that if $\Gamma$ is a $G$-symetric graph
and $\cB$ is a non-trivial $G$-invariant partition of the vertices,
then the graph $\Gamma$ ``decomposes''
into the triple $(\Gamma_\cB, \Gamma[B,C], \cD(B))$.

We may ask now, to what extent the combinatorial structure of $\Gamma$
is determined by the triple $(\Gamma_\cB, \Gamma[B,C], \cD(B))$?
Do these three components contain sufficient information
to reconstruct $\Gamma$ uniquely?
If not, what extra information is required?

%The problem then becomes that of reconstructing $\Gamma$ 
%from the triple $(\Gamma_\cB, \Gamma[B,C], \cD(B))$.


Suppose we are given a symmetric graph $\Lambda$,
a symmetric bipartite graph $\Sigma$ and a flag transitive design
$\cD$ without any particular groups acting upon them.
Let us say that a symmetric graph $\Gamma$ is 
``product'' of $(\Lambda, \Sigma, \cD)$ if there is a group $G$
acting $\Gamma$ and a $G$-invariant partition $\cB$ of
the vertices of $\Gamma$
such that the triple $\Gamma$ decomposes into the triple
$(\Lambda, \Sigma, \cD)$ 

What are the necessary and sufficient
conditions under which these three components $(\Lambda, \Sigma, \cD)$
can be ``glued together'' into some ``product'' $\Gamma$ ?

If the necessary conditions are satisfied,
could there be more than one way to ``glue together'' a given 
triple $(\Lambda, \Sigma, \cD)$ into an imprimitive symmetric graph 
$\Gamma$ ?

%\section{Rephrasing the Problem}
%
%If we ignore the induced bipartite graph for a moment,
%then the map $\Gamma \mapsto (\Gamma_\cB, \cD(B))$
%actually induces a functor from $G$-Graph into the caretesian product
%category ($G$-Flag, $G$-Graph). 
%Well, to be rigerous we've got to describe what happens to the morphisms,
%but its the obvious thing.
%Let us call this functor $D$.
%It is not hard to find examples to show that this functor is not injective.
%
%In the next section we shall define a new category,
%called $G$-Designs, and show that there exists a forgetful functor $F$
%from $G$-Graph into $G$-Design.
%We will then define a functor $E$ from $G$-Design to ($G$-Design, $G$-Flag)
%analogous to $D$.
%We will show that if $F'$ is the functor from ($G$-Graph, $G$-Flag) to 
%($G$-Design, $G$-Flag)
%given by $F'(x,y) = (F(x),y)$ then $E$
%has the property that $E \circ F(x) = F' \circ D(x)$ for any $x \in G$-graph.
%Finally we will conject that $E$ is in fact injective.

%\newpage
%\section{This goes elsewhere}
%We must pause for a moment to highlight an important distinction
%between a symmetric graph and a $G$-symmetric graph.
%We say that a graph $\Gamma$ is symmetric if it is $\Aut(\Gamma)$-symmetric.
%Simmilarly we say that a design $\cD$ is flag-transitive if it is 
%$\Aut(\cD)$-flag transitive.
%
%It is often the case however that for a given symmetric graph $\Gamma$
%there is a subgroup $G$ of the full automorphism group
%which also acts symmetrically on the graph.
%In this case the graph will be $G$-symmetric as well as $\Aut(\Gamma)$-symmetric.
%We saw an example of this in Chapter 1 when we constructed the cube
%as an orbital graph of $(S_4, \cos(C_3))$ but found its full automorphism group
%is in fact $S_4 \times Z_2$.
%A simmilar ``problem'', also arises with flag-transitive designs.
%
%For a given symmetric graph $\Gamma$ let $\Grp(\Gamma)$ denote
%the set of subgroups (upto isomorphism) of $\Aut(\Gamma)$ which act symmetrically
%on $\Gamma$. Simmilarly, for a given flag-transitive design, $\cD$
%let $\Grp(\cD)$ denote the set of subgroups (upto isomorphism) 
%of $\Aut(\cD)$ which act symmetrically on $\cD$.
%Then, for a given group $G$, $\Gamma$ is $G$-symmetric 
%if and only if $G$ admits some element of $\Grp(\Gamma)$ as a composition factor. 
%Simmilarly $\cD$ is $G$-flag transitive if and only if
%$G$ admits some element of $\Grp(\cD)$ as a composition factor.
%
%The study of $G$-symmetric graphs can thus be approached in two different ways.
%The first is to fix the group $G$ and examine the family of graphs $\Gamma$
%which admit a symmetric action of $G$, in the hope that this will reveal
%some information about the group $G$.
%The second approach is to fix a graph $\Gamma$ and study the groups
%which act symmetrically on $\Gamma$ in the hope that this will
%reveals some information about the graph $\Gamma$.

%\section{Also to go elsewhere}

%In chapter 1, we saw that there were two different orbital
%graphs associated with the permutation group $(S_4, \cos(C3))$.
%which admitted the tetrahedron as a quotient.
%One of these was isomorphic to the cube and the other was
%isomorphic to two vertex disjoint tetrahedra.
%One can check that in both cases ``induced bipartite graph''
%of the decomposition
%is isomorphic to $2K_2$ and the ``cross-sectional geometry''
%is the highly degenerate $(2,2,3)$-design consisting of the same
%$2$-element block repeated $3$ times.

%Thus we know that the answer to this last question is definitely
%``yes''. For a given triple $(\Lambda, \Sigma, \cD)$
%there may be more than one way of ``gluing'' them together
%into some imprimitive symmetric graph $\Gamma$.
%The triple $(\Gamma_\cB, \Gamma[B,C], \cD)$ does not, in general,
%contain sufficient information to reconstruct $\Gamma$ uniquely.
